% ==========================================================================
% Anexos
% Estado: se registra en ../STATUS.md, no aquí.
% Se incluye después de \appendix en main.tex, así que cada \chapter aquí
% numera como "Anexo A", "Anexo B", etc. — usar \chapter, no \section.
%
% Debe contener (project-context/formato-anteproyecto.md, "Anexos"):
%   - Análisis de participación (involucrados).
%   - Árbol de problemas (soporta 02-descripcion-problema.tex).
%   - Árbol de objetivos (soporta 04-objetivos.tex).
%   - Tablas, figuras, diagramas o fragmentos de código que den claridad a
%     ideas del cuerpo del documento.
% ==========================================================================

% [verify: ADR-010 — pendiente: el anexo de análisis de participación]

\chapter{Árbol de problemas}\label{anx:arbol-problemas}

Este anexo presenta el árbol de problemas que respalda la descripción del problema del \cref{cha:descripcion-problema}. El árbol reúne en una sola vista la causa principal, los factores causales, el problema central y los efectos que se derivan de él, de modo que el lector pueda comprobar la jerarquía que el capítulo describe en prosa.

La \cref{fig:arbol-problemas} se lee de abajo hacia arriba, y las flechas indican la dirección de la relación causal. En la base se ubica la causa principal, de la que se desprenden tres factores causales que, en conjunto, originan el problema central. De ese problema se derivan dos efectos principales, que a su vez producen efectos secundarios agrupados según quién los experimenta, ya sea la comunidad académica, el personal administrador o la propia infraestructura.

\begin{figure}[htbp]
  \centering
  \hyphenpenalty=10000 \exhyphenpenalty=10000
  \begin{tikzpicture}[
      font=\footnotesize,
      box/.style={draw, rounded corners=2pt, align=center, inner sep=4pt, text width=3.35cm},
      eff/.style={box, fill=black!6},
      cau/.style={box, fill=white},
      prob/.style={box, fill=black!22, very thick, text width=10.6cm},
      wide/.style={box, text width=4.9cm},
      rowlabel/.style={font=\scriptsize\bfseries, rotate=90, anchor=center, align=center},
      flow/.style={-{Latex[length=2.2mm]}, thick}
    ]
    % --- Efectos secundarios ---------------------------------------------
    \node[eff, minimum height=4.2cm] (s1) at (-3.85,0.2)  {\textbf{Estudiantes y profesores}\\[2pt] Validación práctica limitada de sus investigaciones o costos en plataformas comerciales de nube};
    \node[eff, minimum height=4.2cm] (s2) at (0,0.2)   {\textbf{Administradores}\\[2pt] Sobrecarga operativa recurrente por la gestión individual de credenciales y solicitudes};
    \node[eff, minimum height=4.2cm] (s3) at (3.85,0.2)   {\textbf{Infraestructura compartida}\\[2pt] Riesgo de saturación desordenada de las GPU y de fallas no detectadas que pueden degradar la disponibilidad de los nodos};
    % --- Efectos principales ---------------------------------------------
    \node[wide, fill=black!12, minimum height=1.8cm] (e1) at (-2.8,-4.4) {Subutilización sostenida de una infraestructura en la que la universidad ya invirtió};
    \node[wide, fill=black!12, minimum height=1.8cm] (e2) at (2.8,-4.4)  {Interrupción del ciclo de vida de los modelos tras la etapa de entrenamiento};
    % --- Problema central ------------------------------------------------
    \node[prob, minimum height=1.9cm] (p) at (0,-7.9) {\textbf{Problema central}\\[2pt] La infraestructura de cómputo del IAsLab carece de mecanismos para operar como un entorno multiusuario compartido, gobernado y observable};
    % --- Factores causales (causas secundarias) ---------------------------
    \node[cau, minimum height=3.7cm] (c1) at (-3.85,-11.9) {Ausencia de un plano de servicio que gestione motores de inferencia, por lo que el despliegue se realiza a mano sobre el sistema operativo anfitrión};
    \node[cau, minimum height=3.7cm] (c2) at (0,-11.9)    {Falta de integración entre el proveedor de identidad (SAAMFI) y políticas de cuotas, por lo que los roles no se traducen en cupos y los accesos se asignan a mano};
    \node[cau, minimum height=3.7cm] (c3) at (3.85,-11.9)  {Inexistencia de un sistema de observabilidad centralizado, por lo que no hay visibilidad en tiempo real del estado de los nodos ni de la carga de inferencia que atienden};
    % --- Causa principal -------------------------------------------------
    \node[prob, fill=white, minimum height=1.6cm] (r) at (0,-15.4) {\textbf{Causa principal}\\[2pt] Ausencia de una capa de \emph{software} de orquestación y gestión de la etapa de inferencia sobre la infraestructura de cómputo del laboratorio};

    % --- Relaciones causales (de abajo hacia arriba) ------------------------
    \foreach \c in {c1,c2,c3} {
      \draw[flow] (r.north -| \c.south) -- (\c.south);
      \draw[flow] (\c.north) -- (\c.north |- p.south);
    }
    \draw[flow] (p.north -| e1.south) -- (e1.south);
    \draw[flow] (p.north -| e2.south) -- (e2.south);
    \coordinate (bus) at (0,-2.9);
    \draw[thick] (e1.north) -- (e1.north |- bus);
    \draw[thick] (e2.north) -- (e2.north |- bus);
    \draw[thick] (s1.south |- bus) -- (s3.south |- bus);
    \foreach \s in {s1,s2,s3} {
      \draw[flow] (\s.south |- bus) -- (\s.south);
    }

    % --- Etiquetas de nivel --------------------------------------------------
    \node[rowlabel] at (-6.0,0.2)   {Efectos secundarios};
    \node[rowlabel] at (-6.0,-4.4)  {Efectos principales};
    \node[rowlabel] at (-6.0,-7.9)  {Problema};
    \node[rowlabel] at (-6.0,-11.9) {Causas secundarias};
    \node[rowlabel] at (-6.0,-15.4) {Causa principal};
  \end{tikzpicture}
  \caption[Árbol de problemas]{Árbol de problemas del proyecto. Las flechas indican la relación causal y el diagrama se lee de abajo hacia arriba. Elaboración propia a partir del \cref{cha:descripcion-problema}.}
  \label{fig:arbol-problemas}
\end{figure}

\chapter{Árbol de objetivos}\label{anx:arbol-objetivos}

Este anexo presenta el árbol de objetivos que respalda la formulación de los objetivos del \cref{cha:objetivos}. Se obtuvo al plantear en positivo cada elemento del árbol de problemas del Anexo~\labelcref{anx:arbol-problemas}, de modo que cada factor causal se convierte en un medio, el problema central en el objetivo central y los efectos en fines.

La \cref{fig:arbol-objetivos} se lee de abajo hacia arriba. Cada medio conduce a uno de los tres objetivos específicos de construcción, y el conjunto de medios conduce al objetivo general, del que se derivan los fines principales y los fines secundarios. Las columnas conservan el orden del Anexo~\labelcref{anx:arbol-problemas}, de modo que cada medio queda frente a la causa que atiende. Los textos de los objetivos específicos dentro de las cajas están resumidos, y su redacción completa se encuentra en la \cref{sec:objetivos-especificos}.

El \cref{obj:evaluacion} no aparece en el árbol porque ninguna causa del árbol de problemas se asocia a él. Ese objetivo verifica el desempeño de la plataforma una vez construida.

\begin{landscape}
\begin{figure}[p]
  \centering
  \hyphenpenalty=10000 \exhyphenpenalty=10000
  \begin{tikzpicture}[
      font=\footnotesize,
      box/.style={draw, rounded corners=2pt, align=center, inner sep=4pt, text width=5.5cm},
      fin/.style={box, fill=black!6},
      fin2/.style={box, fill=black!12, text width=7.4cm},
      cen/.style={box, fill=black!22, very thick, text width=16.2cm},
      med/.style={box, fill=white},
      rowlabel/.style={font=\scriptsize\bfseries, rotate=90, anchor=center, align=center},
      flow/.style={-{Latex[length=2.2mm]}, thick}
    ]
    % --- Fines secundarios -----------------------------------------------
    \node[fin, minimum height=2.3cm] (a1) at (-5.9,0) {\textbf{Estudiantes y profesores}\\[2pt] Validación práctica de sus investigaciones sobre infraestructura propia, sin costos de plataformas comerciales de nube};
    \node[fin, minimum height=2.3cm] (a2) at (0,0)    {\textbf{Administradores}\\[2pt] Menor sobrecarga operativa en la gestión de credenciales y solicitudes};
    \node[fin, minimum height=2.3cm] (a3) at (5.9,0)  {\textbf{Infraestructura compartida}\\[2pt] Menor riesgo de saturación desordenada de las GPU y de fallas no detectadas, con la disponibilidad de los nodos preservada};
    % --- Fines principales -----------------------------------------------
    \node[fin2, minimum height=1.3cm] (b1) at (-4.1,-3.0) {Aprovechamiento sostenido de una infraestructura en la que la universidad ya invirtió};
    \node[fin2, minimum height=1.3cm] (b2) at (4.1,-3.0)  {Continuidad del ciclo de vida de los modelos tras la etapa de entrenamiento};
    % --- Objetivo central --------------------------------------------------
    \node[cen, minimum height=2.3cm] (c) at (0,-5.9) {\textbf{Objetivo central y objetivo general}\\[2pt] La infraestructura de cómputo del IAsLab opera como un entorno multiusuario compartido, gobernado y observable. Para lograrlo, el proyecto busca construir la plataforma de orquestación y gobernanza de cargas de inteligencia artificial y aprendizaje automático del IAsLab, para consolidar el ciclo de vida de MLOps \emph{on-premise} en la Universidad Icesi};
    % --- Medios y objetivos específicos ---------------------------------------
    \node[med, minimum height=3.6cm] (m1) at (-5.9,-10.1) {\textbf{Plano de servicio que gestiona motores de inferencia}\\[3pt] \Cref{obj:orquestacion}. Implementar el despliegue concurrente de modelos de inteligencia artificial ya entrenados, mediante motores de inferencia empaquetados como imágenes de contenedor};
    \node[med, minimum height=3.6cm] (m2) at (0,-10.1)    {\textbf{Integración entre el proveedor de identidad (SAAMFI) y políticas de cuotas}\\[3pt] \Cref{obj:gobernanza}. Implementar un sistema de cuotas y restricciones acoplado a los roles de SAAMFI, que reserve cupos de cómputo para las cargas de inferencia};
    \node[med, minimum height=3.6cm] (m3) at (5.9,-10.1)  {\textbf{Sistema de observabilidad centralizado}\\[3pt] \Cref{obj:telemetria}. Desarrollar un módulo de observabilidad para capturar y visualizar métricas de \emph{hardware} y de carga de trabajo de los nodos de cómputo};

    % --- Relaciones (de abajo hacia arriba) ------------------------------------
    \foreach \m in {m1,m2,m3} {
      \draw[flow] (\m.north) -- (\m.north |- c.south);
    }
    \draw[flow] (c.north -| b1.south) -- (b1.south);
    \draw[flow] (c.north -| b2.south) -- (b2.south);
    \coordinate (bus) at (0,-1.75);
    \draw[thick] (b1.north) -- (b1.north |- bus);
    \draw[thick] (b2.north) -- (b2.north |- bus);
    \draw[thick] (a1.south |- bus) -- (a3.south |- bus);
    \foreach \a in {a1,a2,a3} {
      \draw[flow] (\a.south |- bus) -- (\a.south);
    }

    % --- Etiquetas de nivel ------------------------------------------------------
    \node[rowlabel] at (-9.35,0)    {Fines secundarios};
    \node[rowlabel] at (-9.35,-3.0) {Fines principales};
    \node[rowlabel] at (-9.35,-5.9) {Objetivo central};
    \node[rowlabel] at (-9.35,-10.1) {Medios y objetivos específicos};
  \end{tikzpicture}
  \caption[Árbol de objetivos]{Árbol de objetivos del proyecto. Las flechas indican que cada medio conduce al objetivo central y que este conduce a los fines, de modo que el diagrama se lee de abajo hacia arriba. Elaboración propia a partir del Anexo~\labelcref{anx:arbol-problemas} y del \cref{cha:objetivos}.}
  \label{fig:arbol-objetivos}
\end{figure}
\end{landscape}
